%
% Halaman Abstrak
%
% @author  Andreas Febrian
% @version 1.00
%
%\chapter*{Abstrak}
\setstretch{1}
\vspace*{0.2cm}
%    \begingroup
    \singlespacing
	\setlength{\parindent}{0pt}
	
	\begin{tabular}{@{}l l p{10cm}}
		Nama&: & \penulis \\
		Program Studi&: & \program \\
		Judul&: & \judul \\
		Pembimbing&: & \pembimbing \\
	\end{tabular}


\bigskip
\bigskip
Autism Spectrum Disorder (ASD) merupakan gangguan perkembangan saraf yang ditandai oleh hambatan dalam komunikasi sosial dan perilaku repetitif. Berbagai penelitian telah memanfaatkan sinyal electroencephalography (EEG) untuk membantu proses klasifikasi ASD, namun keterbatasan jumlah data sering menjadi kendala dalam pengembangan model pembelajaran mesin yang andal. Penelitian ini bertujuan untuk menganalisis efektivitas augmentasi data berbasis Variational Autoencoder (VAE) dalam meningkatkan performa klasifikasi ASD menggunakan sinyal EEG. Penelitian dilakukan dengan membangun model VAE untuk menghasilkan data EEG sintetis yang kemudian digunakan sebagai data augmentasi pada proses klasifikasi. Kualitas data sintetis dievaluasi menggunakan metrik Fréchet Inception Distance (FID), Sliced Wasserstein Distance (SWD), dan normalisasi Euclidean Distance (ED), sedangkan evaluasi klasifikasi dilakukan menggunakan model Deep Neural Network (DNN) dan Extreme Gradient Boosting (XGBoost). Hasil eksperimen menunjukkan bahwa data sintetis yang dihasilkan mampu mempertahankan karakteristik utama sinyal EEG serta memberikan peningkatan performa klasifikasi. Secara signifikan, penggunaan data sintetis pada model XGBoost meningkatkan akurasi dari 97,44\% menjadi 99,54\% sensitivitas dari 97,10\% menjadi 98,67\% spesifisitas dari 97,78\% menjadi 99,96\% dan presisi dari 97,78\% menjadi 99,91\%. Temuan ini menunjukkan bahwa augmentasi berbasis VAE mampu memperkaya variasi data pelatihan dan meningkatkan kemampuan model dalam membedakan kelas ASD dan non-ASD. Penelitian ini menyimpulkan bahwa VAE merupakan pendekatan yang menjanjikan untuk mengatasi keterbatasan data EEG dan meningkatkan performa sistem klasifikasi ASD berbasis EEG.
\newline
\newline
Kata kunci: Autism Spectrum Disorder, Electroencephalography, Augmentasi Data, Variational Autoencoder, Deep Neural Network, XGBoost.

\newpage